\label{sec:results}

\subsection*{Finding 1: Instructed Loadability Replicates Across Models but Scales with Size and Lineage}

Instructed loading---writing a verbalizable concept code into the residual stream---succeeds
only in the workspace band and scales with model size and training objective (Figure~\ref{fig:load}).
Qwen3-4B achieves 10\% hit-rate@5 (gate L1), Qwen3.6-27B 55\% (gate L1b), and the PWM twin 0\% (gate L2).
Uninstructed controls and shuffled nulls separate true signal from byte-fallback artifacts.
The recalibration from union-top-K to model-top-K support was essential: the original metric's structural null floor ($-0.72$)
obscured the true pattern.

\begin{figure}[t]\centering
\includegraphics[width=.72\linewidth]{fig1_loadability}
\caption{Instructed concept loadability in a fixed depth-band across models
(gates: L1, L1b, L2 runs 1--3). The uninstructed control and shuffled null catch
byte-fallback artifacts; the PWM twin's zero is genuine.}\label{fig:load}
\end{figure}

\subsection*{Finding 2: Timing Matters---Sphura\d{t}\d{t}\=a-Gated Writes Beat Fixed Schedules}

Under sampling (the regime where freedom cost is negative), event-gated writes outperform
rate-matched and continuous baselines at matched amplitude (Figure~\ref{fig:dose}).
The dose grid canonical clean-stream re-measurement (gate L18-l8redo, $\alpha \in \{0.02, 0.05, 0.1\}$)
supersedes the pre-stream-fix L8 levels (which ran $\sim 0.1$ higher).
Write \emph{count} is operationally free at this scale (throughput within noise of baseline),
so schedule choice is about behavioral margins, not efficiency.

\begin{figure}[t]\centering
\includegraphics[width=.72\linewidth]{fig2_dose_grid}
\caption{Dose grid on the Qwen3-4B twin under sampling --- \emph{canonical clean-stream
re-measurement} (gate L18-l8redo, superseding the pre-stream-fix L8 levels; the gated
arm ran exactly $0.1$ higher pre-fix, other arms shifted by varying amounts, single
seed). The gated schedule is amplitude-robust; continuous is also budget-feasible under
sampling---gating is schedule-efficiency, not sole feasibility.}\label{fig:dose}
\end{figure}

\subsection*{Finding 3: The Core Claim---Six Independent-Stream Seeds Confirm Gated Steering within Budget}

Gated writes achieve full concept steering within the $\pm 0.5$-nat freedom budget
at six independent-stream seeds (42, 123, 777, 888, 999, 1234), with lift 0.30--0.35.
The alignment advantage over rate-matched controls is small (+0.07...+0.12) but positive in every seed
(one-sided sign test $p\approx0.016$; Figure~\ref{fig:core}, gate L11-rep).
This is the core program claim and sits at confirm tier.

\begin{figure}[t]\centering
\includegraphics[width=.98\linewidth]{fig3_core_claim}
\caption{The core claim at six independent-stream seeds (gates L9-alignconf, L11-rep):
gated lift 0.30--0.35 within the $\pm$0.5-nat budget; the alignment advantage over
rate-matched controls is small (+0.07..+0.12) but positive in every seed
(one-sided sign test $p\approx0.016$).}\label{fig:core}
\end{figure}

\subsection*{Finding 4: Generality via Recipe---Cross-Plant Transfer with Amplitude Calibration}

One-shot transfer with borrowed geometry fails (+0.05 on Qwen3-4B); the calibration recipe
(band-exit lens $\rightarrow$ site probe $\rightarrow$ amplitude scaled to inverse lens transport strength)
reproduces the result: gated $+0.40$ vs prefill $+0.17$ within budget (gate L13).
At four independent-stream seeds (42, 123, 777, 888), the calibrated recipe delivers
gated lift 0.33--0.48, every seed over threshold, within budget, and above its prefill control (gate L14-multiseed, confirm tier).
Qwen3's band-exit Jacobians are $\sim 10 \times$ weaker than Nemotron's; working amplitude must scale inversely.

\begin{figure}[t]\centering
\includegraphics[width=.98\linewidth]{fig4_recipe_and_articulation}
\caption{Left: the per-plant calibration recipe rescues cross-plant transfer (gates L10,
L13). Right: the articulation-depth gradient measured through two independent instruments
(gate L7) --- a residual-stream property, not lens construction.}\label{fig:recipe}
\end{figure}

\subsection*{Finding 5: Amplitude Dose Law---Monotone Lift per Seed on the Weak-Transport Plant}

Sweeping amplitude $\alpha = \mathrm{cap}$ over $\{0.1, 0.2, 0.3, 0.45\}$ on Qwen3-4B,
gated lift rises monotonically \emph{per seed} at three independent-stream seeds
($0 \rightarrow \sim 0.18 \rightarrow 0.4$--$0.48 \rightarrow 0.72$--$0.78$; gates L14-amp, L15-amp-joint),
with no saturation in the tested range, while trajectory-entropy cost stays flat and far inside the $\pm 0.5$-nat budget
(gate L15-amp-joint).
On Nemotron (strong-transport plant), a coarse grid was flat (saturated), prompting a fine-grained sweep
below its saturation point: at $\alpha \in \{0.002, 0.005, 0.01, 0.02\}$ gated lift rises
$0.03 \rightarrow 0.15 \rightarrow 0.28$ and saturates (gate L16-fine, single seed).
Each plant shows within-grid monotone dose response in its own active range ---
Nemotron at $\alpha \sim 0.002$--$0.01$, Qwen3 at $0.1$--$0.45$, $\sim 1.5$ orders of magnitude apart,
consistent with amplitude $\propto 1/$lens-strength (Figure~\ref{fig:scaling}).

\begin{figure}[t]\centering
\includegraphics[width=.98\linewidth]{fig7_scaling_law}
\caption{The amplitude dose law across two plants (gates L14-amp, L15-amp-joint,
L16-fine, L8 replay): per-seed monotone lift on the weak-transport plant (three seeds)
and, on a log-$\alpha$ axis, the strong-transport plant's own rise-then-saturate curve
an order of magnitude lower --- amplitude $\propto$ 1/lens-strength (left). Freedom
cost flat and within budget throughout (right).}
\label{fig:scaling}
\end{figure}

\subsection*{Finding 6: System Architecture---The Research Loop as Active-Inference Selection}

The research program itself ran as expected-free-energy selection over registered experiment menus,
with twenty-four cycles, three menus, fourteen adversarial reviews, and every honest negative
in a replayable ledger (Figures~\ref{fig:arch}, \ref{fig:ledger}).
The loop debugged itself: miscalibrated cost model, replay bugs, sampling-seed correlations (the stream fix),
and staleness violations were all caught and corrected on the record.

\begin{figure}[t]\centering
\includegraphics[width=.86\linewidth]{fig6_architecture}
\caption{System architecture: controller--actuator--plant with the auto-research loop and
adversarial review harness around it.}\label{fig:arch}
\end{figure}

\begin{figure}[t]\centering
\includegraphics[width=.8\linewidth]{fig5_compute_ledger}
\caption{Compute ledger: the entire program to date on one DGX Spark (GB10), guard-shared
with co-resident training jobs, contention recorded in every gate.}\label{fig:ledger}
\end{figure}

\subsection*{Finding 7: PWM Integration---Cold CittaStore Recall is Integration-Ready and Comparable to (Not Equivalent to) the Analytic Baseline}

Integration of the Pratyabhijñā World Model's CittaStore write pathway was blocked until session L20 due to PWM stack availability on GB10.
The frozen (untrained) Hopfield-store initialization—populated with no steering patterns—produces write vectors that route through the band-exit Jacobian and steer within the $\pm0.5$-nat budget.
On three independent-stream seeds (42, 123, 777), cold-store lift matches: $+0.4445$ (seed 42), $+0.5556$ (seed 123), $+0.4445$ (seed 777).

Against the analytic J$^T$u baseline:
\begin{itemize}
\item Seed 42: gap $= 0.0000$ (exact match)
\item Seed 123: gap $= 0.0000$ (exact match)
\item Seed 777: gap $= +0.1111$, which \emph{fails} the registered $|\text{gap}|\le0.05$ equivalence criterion (2.2$\times$ the threshold). Two facts both hold: the miss is the smallest the metric can resolve (one concept-hit, $1/9$) and it is in the cold store's favour ($0.4445$ vs.\ $0.3334$), so it is not a degradation. It is \emph{not} ``within noise''---a nine-sample corpus cannot adjudicate a one-hit difference as signal versus draw.
\end{itemize}

The cold store's \emph{comparability} to the lens-transpose baseline confirms that the PWM write pathway is \emph{plumbing-correct} and dimensionally compatible; it does not establish equivalence, which the registered criterion rejects on one of three seeds. A determinism probe (raised by adversarial review) confirmed the two arms are genuine independent draws---all nine per-generation records differ between arms at seed 123---so the near-mirror aggregate entropy at that seed is a mean-coincidence, not a shared-RNG artifact.
The interesting open question—does training the store against steering-success targets improve over this cold initialization?—remains unconfirmed; the infrastructure to test it is now live.
Tier: confirm (functional cold-recall, 3 seeds, integration verified); fail-on-margin on strict equivalence. The load-bearing result is the integration and the functional in-budget steering; strict equivalence to analytic is not claimed (gate \texttt{gate\_L20\_confirm.json}).

\subsection*{Finding 8: Gated Steering is a Legibility-Efficiency Tradeoff, Not a Raw-Lift Maximizer}

The doctrine's core promise---event-gated writes steer the workspace---holds within entropy budget.
However, a natural question arises: does gated steering unlock \emph{concept surface-rate} (CSR) comparable
to naive continuous flooding? The answer is no, but the tradeoff illuminates the doctrine's purpose.

Across three independent seeds (42, 123, 777) on Qwen3-4B using a pre-registered concept corpus
(site 24, $\alpha=0.3$), continuous flooding achieves CSR $= 0.96$ (concepts surface at 96\% of attempted writes).
Entropy-gated writes and prefill-only writes both achieve CSR $= 0.16$ (gate L21-baselines, three-seed consistent result).
Logit-bias steering (a mechanically-orthogonal baseline with no timing structure) reaches CSR $= 0.40$.
The entropy deltas reveal the tradeoff: continuous flooding pays $+0.073$ nats entropy cost;
gated and prefill pay only $-0.020$ nats (within budget), a $\sim 4 \times$ efficiency gain.

\textbf{Interpretation:} Gated writes trade raw concept-surface throughput for internal-model
transparency and efficiency. A model steered by continuous writes floods the residual stream with
the concept code, leaving high entropy (the model is less certain). A model steered by gated writes
internalizes the concept more narrowly (lower entropy surface-rate) but with tight budget discipline.
This is consistent with the Pratyabhij\~n\=a reading: the doctrine is not about forcing concepts to
surface (that was never the goal), but about \emph{reflexive re-cognition}---the model's own workspace
acknowledging the written pattern. Gating enforces that the model \emph{commits} to the suggestion
only at high-confidence moments, a reading of \emph{sphura\d{t}\d{t}\=a} (the conscious flash) as genuine
commitment, not mere presence in the stream.

The pre-registered criterion (gated CSR $\geq 0.2$) is not met; the gate verdict is fail-on-margin
(gate \texttt{gate\_L21\_baselines\_seed\{42,123,777\}}).
This is an honest negative that refines, not negates, the doctrine's scope.

\subsection*{Finding 9: Naive Contrastive Alignment Steering Does Not Improve Attack Success or Truthfulness on an Aligned Model}

A secondary question: does the doctrine's timing mechanism help with downstream alignment tasks
(refusal, truthfulness) via contrastive direction extraction? Pratyabhij\~n\=a does not claim this---the doctrine
is about workspace transparency, not alignment gain. However, because contrastive activation steering
(CAA-style; Rimsky et al., 2024) has shown promise in recent work, we test whether our infrastructure
enables it.

\textbf{Jailbreak robustness (AdvBench):} We extract a contrastive refusal direction from 10 human-labeled
harmful/refusal-instruction pairs, then steer the model's activations in the band using this direction.
Baseline ASR (attack success rate) on AdvBench: 0.25 (Qwen3-4B refuses 75\% of AdvBench by default, an
already-aligned model). Steered ASR: 0.25 (no change, gate L21-jailbreak-seed42).
Interpretation: the model's refusal heuristics are already surface-level and robust to the activation
steering we apply; improving further would require either (1) deeper model-intrinsic refusal mechanisms
(outside the band's scope), or (2) stronger steering methods (e.g., supervised fine-tuning, not steering).
This is an honest negative: we do not claim refusal improvement.

\textbf{Truthfulness (TruthfulQA):} We extract a contrastive truthful-direction from 15 true/false QA pairs
and steer toward truthfulness. Baseline truthfulness-proxy (string-overlap match to reference answer):
$0.697$ over the evaluated subset. Steered truthfulness-proxy: $0.680$ (a slight \emph{decrease}).
Interpretation: the proxy is weak (semantic scoring via LLM-judge would strengthen claims), and more
importantly, a single contrastive direction extracted from a small sample (15 pairs) may not generalize
to the broader truthfulness distribution. Truthfulness is a multi-faceted property; steering in one
activation subspace is unlikely to improve it uniformly. This honest negative (gate L21-truthful-seed42)
should not surprise: alignment at the activation level is not a solved problem, and the interpretability
toolkit (steering, SAEs, circuits) remains focused on transparency, not alignment assurance.

Both experiments (jailbreak and truthfulness) are single-seed screening results with caveats noted in the
gate files. The key finding is negative: band-targeted steering via contrastive directions does not trivially
improve these benchmarks. This is valuable in two ways: (1) it prevents overclaim---steering is not a shortcut
to alignment, and (2) it clarifies the doctrine's scope---Pratyabhij\~n\=a supplies a theory of workspace structure
and transparency, not alignment gain. Future work combining band-steering with other methods (training,
filtering, ensemble voting) remains open.

% Fig8 STATUS: Fire-Case Slice Visualization (PREPARED, GPU-dependent)
% Searching outputs/ for pre-rendered slice PNG... [handled by make_figures.py, Task 5]
% If PNG found: included via \includegraphics{fig8_fire_case_slice}.
% If not found: document in Task 5 Step 3 with GPU dispatch command for future execution.
